

\begin{theorem}[Spectral normality gives a pure trace-preserving gauge]
    \label{thm:cpsv_normal_tp_gauge}
    \lean{MPSTensor.IsNormalTensor.exists_tpGauge}
    \leanok
    \uses{def:nt_cpsv}
    Let $A$ be a normal tensor of positive bond dimension in the sense of
    \cite[Section~2.2]{Cirac2016MPDO_arXiv}. There is a positive definite
    matrix $\sigma$ such that the gauge obtained from $\sigma$ is
    trace-preserving, primitive, and irreducible. No scalar rescaling is
    needed.
\end{theorem}

\begin{proof}\leanok
    \uses{lem:perron_gauge_primitive_iff}
    The irreducible Perron--Frobenius theorem gives positive definite right and
    left eigenvectors with positive eigenvalues $r$ and $t$. The spectral
    normalization gives $r=1$. Adjointness with respect to the trace pairing
    gives $r\,\Tr(\sigma\rho)=t\,\Tr(\sigma\rho)$. Since $\sigma$ and $\rho$
    are positive definite, their pairing is nonzero, hence $t=1$. Thus the
    Perron gauge is a pure similarity. Irreducibility and primitivity are
    invariant under this similarity.
\end{proof}

\begin{theorem}[Spectral normality gives eventual block injectivity]
    \label{thm:cpsv_normal_eventually_injective}
    \lean{MPSTensor.IsNormalTensor.isNormal}
    \leanok
    \uses{def:nt_cpsv, def:normal}
    Every positive-bond-dimension normal tensor in the sense of
    \cite[Section~2.2]{Cirac2016MPDO_arXiv} is injective after blocking by some
    finite length.
\end{theorem}

\begin{proof}\leanok
    \uses{thm:cpsv_normal_tp_gauge, thm:normal_tp_primitive_irred}
    Apply Theorem~\ref{thm:cpsv_normal_tp_gauge}. The trace-preserving,
    primitive, irreducible representative is eventually block-injective, and
    gauge invariance transports this conclusion back to the original tensor.
\end{proof}

\begin{theorem}[A coefficient-form BNT is an algebraic BNT]
    \label{thm:cpsv_bnt_is_bnt}
    \lean{MPSTensor.IsCPSVBasisOfNormalTensors.isBNT}
    \leanok
    \uses{def:bnt, def:bnt_cpsv,
        lem:cpsv_bnt_block_dimensions_positive,
        thm:cpsv_normal_eventually_injective}
    If $(A_j)_{j=1}^{g}$ is a coefficient-form basis of normal tensors for
    $A$, then $(A_j)_{j=1}^{g}$ is an algebraic basis of normal tensors for
    $A$. Only this implication from spectral normality to algebraic normality
    is asserted.
\end{theorem}

\begin{proof}\leanok
    \uses{lem:cpsv_bnt_block_dimensions_positive,
        thm:cpsv_normal_eventually_injective}
    Every block has positive bond dimension, so spectral normality implies
    eventual block injectivity for each block. The spanning identities and
    eventual linear independence are unchanged.
\end{proof}



\begin{theorem}[Non-decaying rectangular overlap determines the bond dimension]
    \label{thm:rect_overlap_norm_one_dim_eq}
    \lean{MPSTensor.dim_eq_of_mpvOverlap_norm_tendsto_one_of_irreducible_TP}
    \leanok
    \uses{def:mpv_overlap, def:irreducible_tensor, def:tp_kraus}
    Let $A$ and $B$ be irreducible trace-preserving tensors of positive,
    possibly different, bond dimensions. If $|O_{AB}(N)|\longrightarrow 1$,
    then their bond dimensions are equal.
\end{theorem}

\begin{proof}\leanok
    \uses{thm:overlap_decay_rect_irr_tp}
    If the bond dimensions differed, the rectangular overlap-decay theorem
    would give $O_{AB}(N)\to0$, and hence $|O_{AB}(N)|\to0$, contradicting the
    assumed limit.
\end{proof}

\begin{theorem}[Phase-equivalent normal tensors are gauge-phase equivalent]
    \label{thm:cpsv_normal_mpv_phase_gauge}
    \lean{MPSTensor.IsNormalTensor.exists_apply_ne_zero}
    \lean{MPSTensor.IsNormalTensor.selfOverlap_tendsto_one}
    \lean{MPSTensor.isNormalTensor_cast_iff}
    \lean{MPSTensor.norm_eq_one_of_gaugePhase_cast_of_isNormalTensor}
    \lean{MPSTensor.MPVBlockPhaseEquiv.dim_eq_and_gaugePhaseEquiv_of_isNormalTensor}
    \leanok
    \uses{def:nt_cpsv, def:mpv_block_phase_equiv,
        def:gauge_phase_equiv}
    Let $A$ and $B$ be normal tensors of positive, possibly different, bond
    dimensions. Suppose that, for some nonzero $\zeta\in\C$,
    $\ket{V^{(N)}(B)}=\zeta^N\ket{V^{(N)}(A)}$ for $N\geq 0$.
    Then the bond dimensions are equal and there is an invertible matrix $X$
    and a scalar $\eta$ of unit modulus such that
    $B^i=\eta X A^iX^{-1}$ for every physical letter $i$. This is the gauge
    recovery used in \cite[Appendix~A]{Cirac2016MPDO_arXiv}.
\end{theorem}

\begin{proof}\leanok
    \uses{thm:cpsv_normal_tp_gauge, thm:rect_overlap_norm_one_dim_eq,
        thm:nt_overlap_norm_one_gauge}
    Apply Theorem~\ref{thm:cpsv_normal_tp_gauge} to both tensors. Gauge
    invariance of matrix product vectors preserves the phase relation, whose
    self-overlap identities imply unit modulus for the phase. The resulting
    rectangular cross-overlap has norm tending to one. If the bond dimensions
    differed, Theorem~\ref{thm:rect_overlap_norm_one_dim_eq} would give a
    contradiction. Thus the bond dimensions agree, and the
    unit-overlap rigidity theorem gives a gauge-phase equivalence between the
    trace-preserving representatives. Conjugating by the two Perron gauges
    gives the asserted relation for $A$ and $B$.
\end{proof}

\begin{theorem}[Normal-tensor overlap dichotomy]
    \label{thm:cpsv_normal_overlap_dichotomy}
    \lean{MPSTensor.IsNormalTensor.overlap_dichotomy}
    \leanok
    \uses{def:nt_cpsv, def:mpv_overlap, def:gauge_phase_equiv}
    Let $A$ and $B$ be normal tensors of positive, possibly different, bond
    dimensions. Their self-overlaps satisfy
    \begin{align}
        O_{AA}(N) & \longrightarrow 1,
        \notag                         \\
        O_{BB}(N) & \longrightarrow 1.
        \notag
    \end{align}
    Moreover, either $O_{AB}(N)\longrightarrow 0$, or their bond dimensions
    are equal; call the common value $D$. Then there are $X\in\GL_D(\C)$ and
    $\zeta\in\C$, with $|\zeta|=1$, such that
    $B^i=\zeta X A^iX^{-1}$ for every physical letter $i$. In the latter case,
    $|O_{AB}(N)|\to 1$. This is
    \cite[Lemma~A.2]{Cirac2016MPDO_arXiv}.
\end{theorem}

\begin{proof}\leanok
    \uses{thm:cpsv_normal_tp_gauge, thm:overlap_decay_rect_irr_tp,
        thm:overlap_decay_irr_tp, thm:gauge_invariance,
        thm:gauge_phase_mpv_scaling, thm:cpsv_normal_mpv_phase_gauge}
    Put both tensors in their pure trace-preserving gauges. If their bond
    dimensions differ, rectangular overlap decay gives $O_{AB}(N)\to0$. If
    their dimensions agree but they are not gauge-phase equivalent, the
    equal-dimension overlap-decay theorem gives the same conclusion. In the
    remaining case, $B^i=\zeta X A^iX^{-1}$. By
    Theorem~\ref{thm:cpsv_normal_mpv_phase_gauge}, $|\zeta|=1$. Hence
    \begin{align}
        O_{AB}(N)
         & =\overline{\zeta}^{\,N}O_{AA}(N),
        \notag                               \\
        |O_{AB}(N)|
         & =|O_{AA}(N)|\longrightarrow1.
        \notag
    \end{align}
\end{proof}

\begin{theorem}[Exact phase alternative for normal tensors]
    \label{thm:cpsv_normal_mpv_phase_alternative}
    \lean{MPSTensor.IsNormalTensor.mpv_phase_alternative}
    \leanok
    \uses{def:nt_cpsv, def:mpv_overlap, def:mpv_state}
    Let $A$ and $B$ be normal tensors of positive, possibly different, bond
    dimensions. Then either $O_{AB}(N)\to0$, or there is a phase
    $\zeta\in\C$, with $|\zeta|=1$, such that
    $\ket{V^{(N)}(B)}=\zeta^N\ket{V^{(N)}(A)}$ for every chain length $N$.
    This is \cite[Corollary~A.3]{Cirac2016MPDO_arXiv}.
\end{theorem}

\begin{proof}\leanok
    \uses{thm:cpsv_normal_overlap_dichotomy,
        thm:gauge_phase_mpv_scaling}
    Apply Theorem~\ref{thm:cpsv_normal_overlap_dichotomy}. In its second
    alternative, gauge-phase covariance gives
    $V^{(N)}(B)_\sigma=\zeta^N V^{(N)}(A)_\sigma$ for every configuration
    $\sigma$, and hence the asserted equality of vectors.
\end{proof}

\begin{theorem}[Unit overlap from eventual proportionality]
    \label{thm:eventual_proportional_normalized_overlap_norm_one}
    \lean{MPSTensor.mpvOverlap_norm_tendsto_one_of_eventually_proportionalMPV₂}
    \leanok
    \uses{def:mpv_overlap}
    Let $A$ and $B$ be tensors, possibly with different bond dimensions, whose
    self-overlaps tend to one. Suppose that, for all sufficiently large $N$,
    there is $c_N\in\C$ such that
    $\ket{V^{(N)}(A)}=c_N\ket{V^{(N)}(B)}$. Then $|O_{AB}(N)|\to1$.
\end{theorem}

\begin{proof}\leanok
    The proportionality relation gives
    \begin{align}
        O_{AB}(N) & =c_NO_{BB}(N),
        \notag                         \\
        O_{AA}(N) & =|c_N|^2O_{BB}(N).
        \notag
    \end{align}
    Hence $|c_N|^2\to1$ and therefore $|c_N|\to1$. The first identity now
    gives $|O_{AB}(N)|\to1$.
\end{proof}

\begin{theorem}[Geometric proportionality phase for normal tensors]
    \label{thm:normal_proportional_mpv_geometric_phase}
    \lean{MPSTensor.EventuallyNonzeroProportionalMPV₂.%
        exists_unit_phase_power_of_isNormalTensor}
    \leanok
    \uses{def:nt_cpsv, def:mpv_state,
        def:eventually_nonzero_proportional_mpv}
    Let $A$ and $B$ be normal tensors of positive, possibly different, bond
    dimensions. Suppose that, for all sufficiently large $N$, their matrix
    product vectors are proportional by a nonzero scalar. Then there is
    $a\in\C$, with $|a|=1$, such that
    $\ket{V^{(N)}(A)}=a^N\ket{V^{(N)}(B)}$ for every positive $N$.
    This is the one-normal-block consequence of
    \cite[Corollary~A.3 and the proof of Theorem~2.10]{Cirac2016MPDO_arXiv}.
    % Maintainer note: this determines only the unit-phase part for normalized
    % normal blocks. The distinct positive-rescaling problem is recorded in
    % docs/paper-gaps/cpsv16_commuting_form_to_sal.tex.
\end{theorem}

\begin{proof}\leanok
    \uses{thm:eventual_proportional_normalized_overlap_norm_one,
        thm:cpsv_normal_mpv_phase_alternative,
        thm:cpsv_normal_mpv_phase_gauge}
    Normality gives $O_{AA}(N)\to1$ and $O_{BB}(N)\to1$, so
    Theorem~\ref{thm:eventual_proportional_normalized_overlap_norm_one} gives
    $|O_{AB}(N)|\to1$. The first alternative in
    Theorem~\ref{thm:cpsv_normal_mpv_phase_alternative} gives
    $O_{AB}(N)\to0$, hence $|O_{AB}(N)|\to0$, a contradiction. In the
    remaining alternative,
    $\ket{V^{(N)}(B)}=\zeta^N\ket{V^{(N)}(A)}$ with $|\zeta|=1$; take
    $a=\zeta^{-1}$.
\end{proof}



\begin{theorem}[Separated normal tensors are eventually independent]
    \label{thm:cpsv_normal_separated_eventually_li}
    \lean{MPSTensor.exists_eventually_linearIndependent_of_normalTensor_blocks_not_gaugePhaseEquiv}
    \leanok
    \uses{def:nt_cpsv, def:blocks_not_gauge_phase_equiv}
    Let $(A_j)_{j=1}^g$ be a finite family of positive-bond-dimension normal
    tensors, no two of which are gauge-phase equivalent after identifying
    equal bond dimensions. Then the matrix product vectors
    $\{\ket{V^{(N)}(A_j)}\}_{j=1}^g$ are linearly independent for all
    sufficiently large $N$, as in
    \cite[Appendix~A]{Cirac2016MPDO_arXiv}.
\end{theorem}

\begin{proof}\leanok
    \uses{thm:cpsv_normal_tp_gauge,
        lem:bnt_finite_index_overlap_orthonormal_li}
    Put every tensor in the pure trace-preserving gauge of
    Theorem~\ref{thm:cpsv_normal_tp_gauge}. Distinct gauge-phase classes have
    overlaps tending to zero, while each self-overlap tends to one. The finite
    Gram matrix therefore tends to the identity and is invertible for all
    sufficiently large lengths. Gauge invariance of matrix product vectors
    gives the conclusion for the original tensors.
\end{proof}
